% RESULTADO_RECUPERADO. II REV10, 07_elipse.tex:1--61,
% 07b_geometria_elipse.tex:73--115 y 12_dualidad_t.tex:27--107.
% Explicitaciones algebraicas: cota H5<2 para verificar la carta y
% sustitución Rstar^4=(499 alpha-eta_ret)/(501 alpha+eta_ret).
\section{Elipse de acción y determinación del radio normalizado}
\label{iv:ellipse:section}

El lector angular y la sección de acción conservan el mismo
antecedente APP--TRIT--TPK. La forma que se construye a
continuación utiliza las representaciones
\eqref{iv:angular:representantes} y una sección positiva
\(\hbar=\hbar_s\), con
\(s\in\{\mathrm{pre},\mathrm{ret}\}\), de
\eqref{iv:action:hbar}. La razón angular determinará la forma;
la sección de acción conservará su escala común.

\subsection{Dominio y semiejes etiquetados}
\label{iv:ellipse:dominio}

Para abreviar, sean \(A=A_{\deg}\) y \(C^*=C^*_{\deg}\).
En la coordenatización \(A\ne0\), \(\hbar>0\) y \(|C^*/A|<1\),
definimos
\begin{equation}
 \xi=\frac{C^*}{A},\qquad
 \eta_{\mathrm{el}}=\operatorname{artanh}\xi,\qquad
 a_S=\sqrt{2\hbar}\,e^{\eta_{\mathrm{el}}/2},\qquad
 b_S=\sqrt{2\hbar}\,e^{-\eta_{\mathrm{el}}/2}.
 \label{iv:ellipse:semiejes}
\end{equation}
Los semiejes y las coordenadas equilibradas \(Q,P\) tienen
unidad de raíz de acción. Los ejes permanecen etiquetados:
el signo de \(\eta_{\mathrm{el}}\) determina cuál es el
mayor. La curva y su producto de semiejes son
\begin{equation}
 \gamma(\theta)=(a_S\cos\theta,b_S\sin\theta),\qquad
 \frac{Q^2}{a_S^2}+\frac{P^2}{b_S^2}=1,\qquad
 a_Sb_S=2\hbar.
 \label{iv:ellipse:curva}
\end{equation}

\begin{proposition}[Pertenencia de las representaciones a la coordenatización elíptica]
\label{iv:ellipse:pertenencia}
La rama positiva de \eqref{iv:angular:representantes}, con los
intervalos \eqref{iv:action:intervalos}, satisface
\(A>0\), \(0<C^*<A\) y \(\hbar_s>0\).
La reversión de la coordenada impar conserva \(|C^*/A|<1\).
\end{proposition}
\begin{proof}
La proposición \ref{iv:action:positividad} da
\(0<\eta_{\mathrm{ret}}<H_5\) y la positividad de ambas
secciones de acción. Para una cota superior de \(H_5\),
\eqref{iv:action:intervalos} implica
\(\alpha<1/125\) y \(\varphi>8/5\). Al conservar sólo
los términos positivos de \eqref{iv:action:h5}, se obtiene
\[
 H_5<
 \frac{\sqrt5}{2}+
 \frac{9}{16\cdot125^2}+
 \frac{7}{48\cdot125^4}.
\]
Como \(5<81/16\), se tiene \(\sqrt5/2<9/8\). Además,
\[
 \frac{9}{16\cdot125^2}+
 \frac{7}{48\cdot125^4}
 =\frac{421882}{11718750000}<\frac1{1000}.
\]
Por tanto, \(H_5<9/8+1/1000<2\). Usando ahora
\(\alpha>7/1000\),
\[
 0<C^*=2(\eta_{\mathrm{ret}}+\alpha)
 <2\left(2+\frac1{125}\right)=\frac{502}{125}
 <7<1000\alpha=A.
\]
El intercambio de canales de
\ref{iv:angular:inversion} cambia \(C^*\) por \(-C^*\)
y conserva \(A\); por ello conserva la desigualdad absoluta.
\end{proof}

\subsection{Fase paramétrica y área orientada}
\label{iv:ellipse:area}

\begin{proposition}[Ley areal y transporte de las coordenadas angulares]
\label{iv:ellipse:ley-areal}
La fase paramétrica \(\theta\) de \(\gamma\) y su ángulo
polar \(\phi\), levantados continuamente desde cero, satisfacen
\begin{equation}
 \phi(\theta)=\arg(a_S\cos\theta+i b_S\sin\theta),\qquad
 \frac{d\phi}{d\theta}
 =\frac{a_Sb_S}{a_S^2\cos^2\theta+b_S^2\sin^2\theta}>0.
 \label{iv:ellipse:fases}
\end{equation}
En una coordenatización donde ambas tangentes estén definidas,
\(\tan\phi=(b_S/a_S)\tan\theta\). El área orientada barrida
desde la fase cero es
\begin{equation}
 \mathcal A(\theta)=
 \frac12\int_0^\theta(QP'-PQ')\,dt=\hbar\theta,
 \qquad
 \mathcal A(1)=\hbar,\quad\mathcal A(2\pi)=h_s.
 \label{iv:ellipse:area-fase}
\end{equation}
\end{proposition}
\begin{proof}
La derivada del argumento de un vector no nulo es
\((QP'-PQ')/(Q^2+P^2)\). Para la curva
\eqref{iv:ellipse:curva}, el numerador vale
\(a_Sb_S=2\hbar\); esto demuestra la derivada y su
positividad. El cociente \(P/Q\) da la relación entre tangentes;
el levantamiento continuo conserva los cuadrantes y las vueltas.
La integral vale \(a_Sb_S\theta/2=\hbar\theta\).
Finalmente, \(h_s=2\pi\hbar_s\) da la identidad de la
vuelta completa.
\end{proof}

Un radián de fase paramétrica barre \(\hbar\). El ángulo
polar conserva su propia derivada en
\eqref{iv:ellipse:fases}. En coordenadas polares, la ley
equivalente es
\[
 d\mathcal A=\frac12\rho_{\partial}(\phi)^2\,d\phi,\qquad
 \rho_{\partial}(\phi)=
 \left(\frac{\cos^2\phi}{a_S^2}
       +\frac{\sin^2\phi}{b_S^2}\right)^{-1/2}.
\]
La fórmula radial se obtiene sustituyendo
\((Q,P)=\rho(\cos\phi,\sin\phi)\) en
\eqref{iv:ellipse:curva}. Con la orientación de
\(\gamma\), la integral de la uno-forma canónica conserva
el signo
\begin{equation}
 \oint_\gamma P\,dQ=-\pi a_Sb_S=-h_s.
 \label{iv:ellipse:accion-orientada}
\end{equation}
En efecto, \(P\,dQ=-a_Sb_S\sin^2\theta\,d\theta\) y
la integral de \(\sin^2\theta\) en una vuelta es \(\pi\).

\subsection{Reciprocidad de forma y orientación simpléctica}
\label{iv:ellipse:reciprocidad}

Definimos las matrices
\begin{equation}
 D(\eta)=\begin{pmatrix}e^\eta&0\\0&e^{-\eta}\end{pmatrix},
 \qquad
 S_2=\begin{pmatrix}0&1\\1&0\end{pmatrix},
 \qquad
 J_2=\begin{pmatrix}0&-1\\1&0\end{pmatrix}.
 \label{iv:ellipse:transportes}
\end{equation}

\begin{proposition}[Forma cuadrática y acción bajo reciprocidad]
\label{iv:ellipse:accion-reciproca}
Ambos transportes satisfacen
\(T^{\mathsf T}D(-\eta)T=D(\eta)\).
Para \(\omega=dQ\wedge dP\), el intercambio \(S_2\) es
antisimpléctico y el giro \(J_2\) es simpléctico:
\[
 S_2^*\omega=-\omega,\qquad J_2^*\omega=\omega.
\]
Si \(\mathcal I(\gamma)=\oint_\gamma P\,dQ\) sobre una
curva cerrada, entonces
\begin{equation}
 \mathcal I(S_2\gamma)=-\mathcal I(\gamma),\qquad
 \mathcal I(J_2\gamma)=\mathcal I(\gamma).
 \label{iv:ellipse:signos-accion}
\end{equation}
\end{proposition}
\begin{proof}
La multiplicación matricial intercambia las dos entradas
diagonales de \(D\), lo que demuestra ambas congruencias.
Los determinantes de \(S_2,J_2\) son \(-1,1\);
de ahí se obtiene la transformación de la forma de área.
Para \(\lambda=P\,dQ\), el cálculo directo da
\[
 S_2^*\lambda=d(PQ)-\lambda,\qquad
 J_2^*\lambda=\lambda-d(PQ).
\]
La integral de la diferencial exacta sobre la curva cerrada
es cero, y quedan probadas las dos identidades de acción.
\end{proof}

La orientación distingue los dos transportes aunque publiquen
la misma forma cuadrática, incluso en el caso circular.
El marcado de los ejes y el signo de la acción pertenecen,
por tanto, a los datos que se conservan al pasar a la
representación reticular.

\Needspace{12\baselineskip}
\subsection{Radio normalizado y matriz de la lectura reticular}
\label{iv:ellipse:radio}

\begin{proposition}[Determinación y recuperabilidad del radio de forma]
\label{iv:ellipse:radio-prop}
La elipse etiquetada determina un único radio positivo
\begin{equation}
 R_\star=\sqrt{\frac{b_S}{a_S}}
 =e^{-\eta_{\mathrm{el}}/2}
 =\left(\frac{A-C^*}{A+C^*}\right)^{1/4},
 \qquad \tau_{\mathrm{el}}=iR_\star^2.
 \label{iv:ellipse:rstar}
\end{equation}
La razón angular se recupera mediante
\begin{equation}
 \frac{C^*}{A}=\frac{1-R_\star^4}{1+R_\star^4},
 \qquad
 R_\star^4=
 \frac{499\alpha-\eta_{\mathrm{ret}}}
      {501\alpha+\eta_{\mathrm{ret}}}.
 \label{iv:ellipse:rstar-compuesto}
\end{equation}
El intercambio de semiejes transforma
\((\eta_{\mathrm{el}},R_\star,\tau_{\mathrm{el}})\) en
\((-\eta_{\mathrm{el}},R_\star^{-1},-1/\tau_{\mathrm{el}})\)
y conserva \(a_Sb_S=2\hbar\).
\end{proposition}
\begin{proof}
La división de los semiejes da \(b_S/a_S=e^{-\eta_{\mathrm{el}}}\).
La aplicación \(R\mapsto R^2\) es biyectiva sobre los reales
positivos, de modo que la raíz positiva es única.
Para \(\xi=C^*/A\in(-1,1)\),
\[
 2\operatorname{artanh}\xi
 =\log\frac{1+\xi}{1-\xi};
\]
por tanto, \(R_\star^4=(1-\xi)/(1+\xi)\).
Despejar \(\xi\) demuestra la primera identidad de
\eqref{iv:ellipse:rstar-compuesto}. Sustituir
\(A=1000\alpha\) y
\(C^*=2(\eta_{\mathrm{ret}}+\alpha)\) da
\[
 \frac{A-C^*}{A+C^*}
 =\frac{998\alpha-2\eta_{\mathrm{ret}}}
        {1002\alpha+2\eta_{\mathrm{ret}}}
 =\frac{499\alpha-\eta_{\mathrm{ret}}}
        {501\alpha+\eta_{\mathrm{ret}}}.
\]
El intercambio invierte la razón de semiejes y cambia el
signo de \(\eta_{\mathrm{el}}\). La identidad
\(-1/(iR_\star^2)=iR_\star^{-2}\) da el transporte
modular. El producto de los semiejes permanece fijo.
\end{proof}

En la rama positiva, \(0<R_\star<1\); la coordenatización de forma
revertida expresa \(R_\star^{-1}>1\). La razón angular
se recupera a partir del radio, mientras los ángulos absolutos
y su historia permanecen en el estado enriquecido.
La forma circular corresponde a \(C^*/A=0\), con
\(R_\star=1\). Cambiar la sección positiva de acción,
manteniendo el par angular, modifica la escala común de
los semiejes y conserva su razón.

\begin{proposition}[Forma reticular determinada por la elipse]
\label{iv:ellipse:matriz-prop}
La matriz de semiejes cuadrados normalizados es
\begin{equation}
 D(\eta_{\mathrm{el}})
 =\frac1{2\hbar}
    \begin{pmatrix}a_S^2&0\\0&b_S^2\end{pmatrix}
 =\begin{pmatrix}R_\star^{-2}&0\\0&R_\star^2\end{pmatrix}.
 \label{iv:ellipse:matriz}
\end{equation}
Su evaluación sobre un par entero produce
\begin{equation}
 E_\star(m,w)=
 \begin{pmatrix}m&w\end{pmatrix}
 D(\eta_{\mathrm{el}})
 \begin{pmatrix}m\\w\end{pmatrix}
 =\frac{m^2}{R_\star^2}+w^2R_\star^2.
 \label{iv:ellipse:energia}
\end{equation}
La ecuación de la elipse multiplicada por \(2\hbar\)
utiliza, en cambio, la matriz inversa \(D(-\eta_{\mathrm{el}})\).
\end{proposition}
\begin{proof}
Los cuadrados de los semiejes de
\eqref{iv:ellipse:semiejes}, divididos por \(2\hbar\),
son \(e^{\eta_{\mathrm{el}}}\) y
\(e^{-\eta_{\mathrm{el}}}\).
Por \eqref{iv:ellipse:rstar}, éstos son
\(R_\star^{-2}\) y \(R_\star^2\).
La multiplicación por el vector entero demuestra
\eqref{iv:ellipse:energia}. Finalmente,
\(2\hbar/a_S^2=e^{-\eta_{\mathrm{el}}}\) y
\(2\hbar/b_S^2=e^{\eta_{\mathrm{el}}}\), lo que
identifica la matriz de la ecuación elíptica.
\end{proof}

El radio \(R_\star\) es una representación adimensional de
la forma etiquetada. El radio polar
\(\rho_{\partial}\) pertenece a la recta de raíz de acción.
La representación por un radio espacial requiere su coordenatización
dimensional posterior; ésta conserva la forma ya determinada
y la independencia de su construcción respecto de una
longitud prescrita.
